% SEGMENT OLP-0343-S01
% Part: lambda-calculus
% Chapter: introduction
% Section: overview

\documentclass[../../../include/open-logic-section]{subfiles}

\begin{document}

\olfileid{lam}{int}{ovr}
\olsection{மேலோட்டம்}

லாம்டா கலனத்தை அலோன்சோ சர்ச் 1930-களின் தொடக்கத்தில் ஆக்கமுறைத் தருக்கத்திற்கு
ஓர் அடிப்படையாகவே முதலில் வடிவமைத்தார்; கணிக்கத்தக்க சார்புகளின் மாதிரியாக
\emph{அல்ல}. ஆனால் டியூரிங்-கணிக்கத்தக்க சார்புகள், பகுதி மீள்வரையறைச்
சார்புகள் போன்ற கணிப்புத்தன்மையின் பிற வரையறைகளுக்கு அது சமானமானது என்று
விரைவிலேயே காட்டப்பட்டது. இது முதலில் ஓரளவு வியப்பை ஏற்படுத்தியது என்ற
உண்மை, அவ்வகைப்பாட்டை இன்னும் சுவையானதாக்குகிறது.

ஒரு சார்புக்கு ஒரு பெயரை வரையறுப்பதற்குப் பதிலாக, அதை வரையறுக்கும் குறியீட்டுக்
கோவையால் அச்சார்பை நேரடியாகக் குறிப்பிடுவதற்கு லாம்டா குறிமானம் ஒரு வசதியான
வழி. ``$f(x) = x + 3$ ஆல் வரையறுக்கப்படும் சார்பு $f$ ஆகட்டும்'' என்று
சொல்வதற்குப் பதிலாக, ``$f$ என்பது $\lambd[x][(x + 3)]$ என்ற சார்பாகட்டும்''
என்று சொல்லலாம். வேறு சொற்களில், $\lambd[x][(x+3)]$ என்பது தனது மாறியுடன்
மூன்றைக் கூட்டும் சார்பிற்கான ஒரு \emph{பெயர்} மட்டுமே. இக்கோவையில் $x$ ஒரு
துணை மாறி, அதாவது ஓர் இடநிரப்பு: அதே சார்பை $\lambd[y][(y + 3)]$ என்றும்
குறிக்கலாம். வேறு அளவுருக்கள் அருகில் இருந்தாலும் இக்குறிமானம் செயல்படும்.
எடுத்துக்காட்டாக, $g(x, y)$ இரு மாறிகளைக் கொண்ட சார்பாகவும், $k$ ஓர் இயல்
எண்ணாகவும் இருக்கட்டும். அப்போது $\lambd[x][g(x,k)]$ என்பது எந்த~$x$ ஐயும்
~$g(x, k)$ க்கு அனுப்பும் சார்பு.

ஒரு குறியீட்டுக் கோவையிலிருந்து இவ்வாறு சார்பை வரையறுப்பது \emph{லாம்டா
அருவமாக்கம்} எனப்படுகிறது. லாம்டா அருவமாக்கத்தின் மறுபக்கம்
\emph{செயல்படுத்தல்}: ஒரு சார்பு $f$ நம்மிடம் இருந்தால் (எடுத்துக்காட்டாக,
இயல் எண்கள்மீது வரையறுக்கப்பட்டதாக), அதை 2 போன்ற எந்த மதிப்பிலும்
\emph{செயல்படுத்தலாம்}. வழக்கமான குறிமானத்தில் அதன் விளைவாக~$f(2)$ என்று
எழுதுகிறோம்.

லாம்டா அருவமாக்கத்தையும் செயல்படுத்தலையும் சேர்த்தால் என்ன நிகழும்? அருவமாக்கிய
மாறிக்குப் பதிலாகச் செயல்படுத்தப்படும் உறுப்பை ``இடுவதன்'' மூலம் விளையும்
கோவையை எளிமைப்படுத்தலாம். எடுத்துக்காட்டாக,
\[
(\lambd[x][(x + 3)])(2)
\]
என்பதை~$2 + 3$ என எளிமைப்படுத்தலாம்.

இதுவரை வழக்கமான கருத்துகளுக்குப் புதிய குறிமானங்களை அறிமுகப்படுத்தியதைத் தவிர
வேறொன்றும் செய்யவில்லை. எனினும், லாம்டா கலனம் கணக்கோட்பாட்டுப் பார்வையிலிருந்து
மேலும் அடிப்படையான விலகலைக் குறிக்கிறது. இந்தக் கட்டமைப்பில்:
\begin{enumerate}
\item எல்லாமே ஒரு சார்பைக் குறிக்கிறது.
\item லாம்டா அருவமாக்கத்தைப் பயன்படுத்திச் சார்புகளை வரையறுக்கலாம்.
\item எதையும் வேறெதன்மீதும் செயல்படுத்தலாம்.
\end{enumerate}
எடுத்துக்காட்டாக, $F$ லாம்டா கலனத்தின் ஒரு சொல்லாக இருந்தால், $F(F)$ எப்போதும்
பொருளுடையதாகவே கொள்ளப்படுகிறது. இத்தாராளமான கட்டமைப்பு \emph{வகையிடப்படாத}
லாம்டா கலனம் எனப்படுகிறது; இங்கு ``வகையிடப்படாத'' என்பது ``எதன்மீது எதைச்
செயல்படுத்தலாம் என்பதற்குக் கட்டுப்பாடு இல்லை'' என்று பொருள்படும்.

\begin{digress}
வகையிடப்படாத வடிவத்தின் ஒரு முக்கிய மாறுபாடாக \emph{வகையிட்ட} லாம்டா கலனமும்
உள்ளது. பல வழிகளில் வகையிட்ட லாம்டா கலனம் வகையிடப்படாத கலனத்தை ஒத்திருந்தாலும்,
அதைச் செவ்வியல் கணக்கோட்பாட்டுக் கட்டமைப்போடு இசைவிக்க மிகவும் எளிது; மேலும்
அதன் சில பண்புகள் முற்றிலும் வேறுபட்டவை.

லாம்டா கலனத்தைப் பற்றிய ஆராய்ச்சி கோட்பாட்டுக் கணினியியலிலும் நிரலாக்க
மொழிகளின் வடிவமைப்பிலும் மையமானதாக அமைந்துள்ளது. 1950-களில் ஜான் மெக்கார்த்தி
வடிவமைத்த LISP, இக்கருத்துகளால் தாக்கம் பெற்ற ஒரு தொடக்ககால மொழி.
\end{digress}

\end{document}
